% Source: arXiv:2203.12563v3, REsubmission.tex, lines 1695--1777.
% Auxiliary normalized operator identification.

\subsection{The common core for every ordered pair of blocks}
\label{sec:mpo_joint_core_hamiltonian}

Write $A_x=A_x^0$ and $D_x=D_x^0$ in this subsection. The exterior
multiplicities $E_x,F_x$ are arbitrary nonnegative integers. The shared
physical alphabet remains $\Fin{d_0}$ throughout the interior.

\begin{theorem}[Dependent orthogonal projections]
    \label{thm:mpo_dependent_orthogonal_projections}
    \lean{ContinuousLinearMap.isSymmetricProjection_sigmaFiberwiseMap,
        ContinuousLinearMap.isSymmetricProjection_dependentRightFiberwiseMap,
        LinearIsometryEquiv.conj_eq_dependentRightFiberwiseMap_of_ker_iff}
    \leanok
    Let $I_q,S_q$ be finite families of finite sets, and let $P_q$ be
    orthogonal projections on $\C^{I_q}$. Then
    $\bigoplus_q(P_q\otimes I_{S_q})$ is an orthogonal projection.
    If $P$ is an orthogonal projection and
    $U:\mathcal H\to\bigoplus_q\C^{I_q\times S_q}$ is unitary with
    $v\in\ker P$ exactly when $(Uv)_{q,s}\in\ker P_q$ for every $q,s$,
    then $UPU^{-1}=\bigoplus_q(P_q\otimes I_{S_q})$.
    Empty label sets and empty spectator fibers are allowed.
\end{theorem}
\begin{proof}
    \leanok
    The inner product is the sum of the fiber inner products, so symmetry
    and idempotence hold on each fiber. The two conjugated projections
    have the same kernel, hence the same range, the orthogonal complement
    of that kernel. Orthogonal projections with the same range are equal.
\end{proof}

\begin{definition}[Original-tensor boundary cores]
    \label{def:mpo_original_boundary_cores}
    \lean{MPSTensor.MPOSymmetry.endpointLeftCoreMap,
        MPSTensor.MPOSymmetry.endpointRightCoreMap,
        MPSTensor.MPOSymmetry.endpointLeftCoreSupportES,
        MPSTensor.MPOSymmetry.endpointRightCoreSupportES,
        MPSTensor.MPOSymmetry.endpointLeftCoreConstraintES,
        MPSTensor.MPOSymmetry.endpointRightCoreConstraintES}
    \leanok
    Define $L_x:\C^{D_x}\to\C^{D_x\times d_0}$ and
    $R_x:\C^{D_x}\to\C^{d_0\times D_x}$ by
    \begin{align}
        (L_xu)_{b,i} & =\sum_c A_x^i[b,c]u_c,  &
        (R_xu)_{i,c} & =\sum_b u_b A_x^i[b,c].
    \end{align}
    Put $S_{L,x}=\operatorname{ran}L_x$,
    $S_{R,x}=\operatorname{ran}R_x$, and let
    $p_{L,x}=\Pi_{S_{L,x}^{\perp}}$,
    $p_{R,x}=\Pi_{S_{R,x}^{\perp}}$.
\end{definition}

\begin{definition}[Dependent exterior regrouping at an edge]
    \label{def:mpo_joint_edge_regrouping}
    \lean{MPSTensor.MPOSymmetry.jointEndpointFirstEdgeSpectatorEquiv,
        MPSTensor.MPOSymmetry.jointEndpointLastEdgeSpectatorEquiv,
        MPSTensor.MPOSymmetry.jointEndpointFirstEdgeSpectatorIsometry,
        MPSTensor.MPOSymmetry.jointEndpointLastEdgeSpectatorIsometry}
    \leanok
    \uses{def:mpo_joint_edge_cores}
    On the first edge, the unitary coordinate regrouping is
    $((x,a,b),i)\mapsto(x,(b,i),a)$. On the last edge, it is
    $(i,(x,c,e))\mapsto(x,(i,c),e)$.
\end{definition}

\begin{theorem}[The one-sided ranges split over exterior fibers]
    \label{thm:mpo_joint_core_edge_fibers}
    \lean{MPSTensor.MPOSymmetry.mem_jointEndpointFirstEdgeCoreSupportES_iff_fibers,
        MPSTensor.MPOSymmetry.mem_jointEndpointLastEdgeCoreSupportES_iff_fibers}
    \leanok
    \uses{def:mpo_joint_edge_cores,def:mpo_original_boundary_cores,
        def:mpo_joint_edge_regrouping}
    Under the respective regroupings,
    \begin{align}
        v\in\operatorname{ran}\Phi_{L,E}
         & \iff v_{x,a}\in S_{L,x}\text{ for every }x,a, \\
        w\in\operatorname{ran}\Phi_{R,E}
         & \iff w_{x,e}\in S_{R,x}\text{ for every }x,e.
    \end{align}
\end{theorem}
\begin{proof}
    \leanok
    In the first identity a rectangular boundary $Y_x$ supplies
    $u_c=Y_x[c,a]$ at each fixed $x,a$. Conversely, choose such a vector
    independently for each $x,a$ and assemble these columns into $Y_x$.
    The last identity uses the rows $u_b=Z_x[e,b]$ in the same way.
    Neither construction selects a label or an exterior value from an
    empty set.
\end{proof}

\begin{theorem}[Exact normalized edge projector identities]
    \label{thm:mpo_joint_core_edge_projectors}
    \lean{MPSTensor.MPOSymmetry.jointEndpointFirstEdgeCoreConstraint_conj_spectators,
        MPSTensor.MPOSymmetry.jointEndpointLastEdgeCoreConstraint_conj_spectators,
        MPSTensor.MPOSymmetry.jointEndpointFirstEdgeCoreConstraint_apply_fiber,
        MPSTensor.MPOSymmetry.jointEndpointLastEdgeCoreConstraint_apply_fiber}
    \leanok
    \uses{thm:mpo_joint_core_edge_fibers,def:mpo_original_boundary_cores}
    The projections onto the orthogonal complements of
    $\operatorname{ran}\Phi_{L,E}$ and $\operatorname{ran}\Phi_{R,E}$
    become respectively
    $\bigoplus_x(p_{L,x}\otimes I_{E_x})$ and
    $\bigoplus_x(p_{R,x}\otimes I_{E_x})$ under the edge regroupings.
\end{theorem}
\begin{proof}
    \leanok
    \uses{thm:mpo_joint_core_edge_fibers,thm:mpo_dependent_orthogonal_projections}
    Their kernels are exactly the ranges in
    Theorem~\ref{thm:mpo_joint_core_edge_fibers}; apply
    Theorem~\ref{thm:mpo_dependent_orthogonal_projections}.
\end{proof}

\begin{definition}[Ordered-pair chain and core coordinates]
    \label{def:mpo_joint_ordered_pair_coordinates}
    \lean{MPSTensor.MPOSymmetry.JointEndpointLeftIndex,
        MPSTensor.MPOSymmetry.JointEndpointRightIndex,
        MPSTensor.MPOSymmetry.JointEndpointChainCfg,
        MPSTensor.MPOSymmetry.JointEndpointCoreCfg,
        MPSTensor.MPOSymmetry.jointEndpointChainSpectatorEquiv,
        MPSTensor.MPOSymmetry.jointEndpointChainSpectatorIsometry,
        MPSTensor.MPOSymmetry.jointEndpointExteriorFiber}
    \leanok
    \uses{def:mpo_joint_edge_cores}
    For $m\geq1$ interior sites, put
    $I_{x,y,m}=\Fin{D_x}\times\operatorname{Cfg}(d_0,m)\times\Fin{D_y}$.
    The normalized chain coordinates and their unitary regrouping are
    \begin{align}
        \mathcal B_L(E)\times\operatorname{Cfg}(d_0,m)\times\mathcal B_R(F)
                                 & \longrightarrow\coprod_{x,y}
        \bigl(I_{x,y,m}\times(\Fin{E_x}\times\Fin{F_y})\bigr),            \\
        ((x,a,b),\sigma,(y,c,e)) & \longmapsto((x,y),(b,\sigma,c),(a,e)).
    \end{align}
    Denote the induced unitary by $U_{E,F,m}$. Its fiber at
    $(x,y),(a,e)$ is given by
    \begin{align}
        v_{x,y;a,e}(b,\sigma,c)=v((x,a,b),\sigma,(y,c,e)).
    \end{align}
    All ordered pairs occur, including $x\ne y$.
\end{definition}

\begin{definition}[First, middle, and last coordinate placements]
    \label{def:mpo_joint_core_placements}
    \lean{MPSTensor.MPOSymmetry.jointEndpointChainFirstEquiv,
        MPSTensor.MPOSymmetry.jointEndpointChainLastEquiv,
        MPSTensor.MPOSymmetry.jointEndpointChainBulkEquiv,
        MPSTensor.MPOSymmetry.jointEndpointCoreFirstEquiv,
        MPSTensor.MPOSymmetry.jointEndpointCoreLastEquiv,
        MPSTensor.MPOSymmetry.jointEndpointCoreBulkEquiv}
    \leanok
    \uses{def:mpo_joint_ordered_pair_coordinates}
    The first placement separates $((x,a,b),\sigma_1)$ from the remaining
    physical word and right boundary. The last separates
    $(\sigma_m,(y,c,e))$ from the left boundary and preceding word.
    The middle placement separates $\sigma$ from both boundaries.
    On $I_{x,y,m}$ the same placements separate $(b,\sigma_1)$,
    $(\sigma_m,c)$, and $\sigma$, respectively, retaining the remaining
    core coordinates as spectators.
\end{definition}

\begin{definition}[The explicit common core and normalized sum]
    \label{def:mpo_joint_common_core_sum}
    \lean{MPSTensor.MPOSymmetry.jointEndpointNormalizedFirstTerm,
        MPSTensor.MPOSymmetry.jointEndpointNormalizedLastTerm,
        MPSTensor.MPOSymmetry.jointEndpointNormalizedBulkTerm,
        MPSTensor.MPOSymmetry.jointEndpointCoreFirstTerm,
        MPSTensor.MPOSymmetry.jointEndpointCoreLastTerm,
        MPSTensor.MPOSymmetry.jointEndpointCoreBulkTerm,
        MPSTensor.MPOSymmetry.jointEndpointCoreHamiltonian,
        MPSTensor.MPOSymmetry.jointEndpointNormalizedSum}
    \leanok
    \uses{def:mpo_original_boundary_cores,def:mpo_joint_core_placements}
    Let $A_{\mathrm{joint}}=\bigoplus_z A_z$ on the original shared
    physical alphabet, with all block weights equal to one. Its two-site
    parent interaction is denoted by $h_{\mathrm{joint}}$. Define
    \begin{align}
        G_{x,y,m} & =p_{L,x}^{\mathrm{first}}+
        \sum_{j=1}^{m-1}h_{\mathrm{joint}}^{(j,j+1)}+
        p_{R,y}^{\mathrm{last}}.\label{eq:mpo_joint_core_sum}
    \end{align}
    \begin{align}
        K_{E,F,m} & =\Pi_{(\operatorname{ran}\Phi_{L,E})^\perp}^{\mathrm{first}}+
        \sum_{j=1}^{m-1}h_{\mathrm{joint}}^{(j,j+1)}+
        \Pi_{(\operatorname{ran}\Phi_{R,F})^\perp}^{\mathrm{last}}.
        \label{eq:mpo_joint_normalized_sum}
    \end{align}
    The middle interaction is joint even when $x\ne y$; it is not a
    parent interaction chosen from one boundary label. The full physical
    alphabet, including any unused directions, remains in every core.
\end{definition}

\begin{theorem}[The shortest supported chain has no middle edge]
    \label{thm:mpo_joint_three_site_middle}
    \lean{MPSTensor.MPOSymmetry.jointEndpointNormalizedBulkTerm_one_eq_zero}
    \leanok
    \uses{def:mpo_joint_common_core_sum}
    At $m=1$, the middle sum vanishes. The total length is three and
    the first and last edges are distinct.
\end{theorem}
\begin{proof}
    \leanok
    No two-site window fits inside a single middle site, so the middle
    sum is indexed by the empty set.
\end{proof}

\begin{theorem}[Termwise identification on every ordered pair]
    \label{thm:mpo_joint_core_term_fibers}
    \lean{MPSTensor.MPOSymmetry.jointEndpointNormalizedFirstTerm_apply_fiber,
        MPSTensor.MPOSymmetry.jointEndpointNormalizedLastTerm_apply_fiber,
        MPSTensor.MPOSymmetry.jointEndpointNormalizedBulkTerm_apply_fiber,
        MPSTensor.MPOSymmetry.jointEndpointNormalizedSum_apply_fiber}
    \leanok
    \uses{def:mpo_joint_common_core_sum,def:mpo_joint_ordered_pair_coordinates}
    For each fixed $(x,y),(a,e)$, the three terms in
    (\ref{eq:mpo_joint_normalized_sum}) act on $v_{x,y;a,e}$ as the
    corresponding three terms in (\ref{eq:mpo_joint_core_sum}). Hence
    $(K_{E,F,m}v)_{x,y;a,e}=G_{x,y,m}v_{x,y;a,e}$.
\end{theorem}
\begin{proof}
    \leanok
    \uses{thm:mpo_joint_core_edge_projectors,def:mpo_joint_core_placements}
    The first edge identity fixes $x,a$ and leaves the right boundary
    unchanged. The last fixes $y,e$ and leaves the left boundary unchanged.
    The middle placement evaluates the same joint open-chain operator on
    $\sigma$ in both coordinate systems. Add these three identities.
\end{proof}

\begin{theorem}[Exact dependent spectator operator identification]
    \label{thm:mpo_joint_core_operator_identity}
    \lean{MPSTensor.MPOSymmetry.jointEndpointNormalizedSum_intertwines_coreSpectators,
        MPSTensor.MPOSymmetry.jointEndpointNormalizedSum_conj_coreSpectators}
    \leanok
    \uses{def:mpo_joint_common_core_sum,def:mpo_joint_ordered_pair_coordinates}
    For arbitrary $E_x,F_x\geq0$ and $m\geq1$,
    \begin{align}
        U_{E,F,m}K_{E,F,m}U_{E,F,m}^{-1}
         & =\bigoplus_{x,y}\bigl(G_{x,y,m}\otimes I_{E_xF_y}\bigr).
        \label{eq:mpo_joint_core_identity}
    \end{align}
    Thus the choices $E=F=D^0$ and $E=F=D^0+D^1$ use the same concrete
    family $G_{x,y,m}$. Empty labels, zero dimensions, and distinct
    boundary labels are included.
\end{theorem}
\begin{proof}
    \leanok
    \uses{thm:mpo_joint_core_term_fibers}
    Evaluate both sides at an arbitrary $(x,y),(b,\sigma,c),(a,e)$ and
    use Theorem~\ref{thm:mpo_joint_core_term_fibers}.
\end{proof}

% Identifying K with the boundary-deformed actual physical Hamiltonian
% remains a separate reducing-frame and compressed-kernel obligation.
% Equation (eq:mpo_joint_core_identity) alone makes no physical gap claim.

\begin{theorem}[Identical gaps for the two exterior multiplicities]
    \label{thm:mpo_joint_normalized_gap_equivalence}
    \lean{MPSTensor.MPOSymmetry.jointEndpointNormalizedSum_norm_gap_iff_enlarged}
    \leanok
    \uses{thm:mpo_joint_core_operator_identity,thm:mpo_joint_boundary_spectator_gap,
        thm:martingale_isometric_conjugation_gap}
    Suppose $D_x^0>0$ for every label, $D_x^1\geq0$, and $\delta\geq0$.
    Set $D=D^0$ and $\widetilde D=D^0+D^1$. For each $m\geq1$,
    \begin{align}
        \delta\|v\|\leq\|K_{D,D,m}v\|
        \quad(v\perp\ker K_{D,D,m})
        \iff
        \delta\|w\|\leq\|K_{\widetilde D,\widetilde D,m}w\|
        \quad(w\perp\ker K_{\widetilde D,\widetilde D,m}).
    \end{align}
    Each side quantifies over its entire indicated kernel complement.
\end{theorem}
\begin{proof}
    \leanok
    \uses{thm:mpo_joint_core_operator_identity,thm:mpo_joint_boundary_spectator_gap,
        thm:martingale_isometric_conjugation_gap}
    The coordinate unitaries preserve norms, kernels and their orthogonal
    complements. Apply (\ref{eq:mpo_joint_core_identity}) on each side and
    then Theorem~\ref{thm:mpo_joint_boundary_spectator_gap} to the common
    family $G_{x,y,m}$. The exterior fibers are nonempty because $D_x^0>0$.
\end{proof}
